\section{The experiment that can separate intelligence from delivery}

The thesis is testable only if the comparison starts from the same task, not from unrelated demonstrations. A task episode should specify an owner-authorized goal, the initial state of relevant services, the permitted routes, a deadline, and an external acceptance check. The dataset should include tasks that stock products usually complete, tasks that produce an observed stop, and tasks that should end in a correct denial or human handoff. A sample composed only of failures selected after seeing a stock product decline would exaggerate the addressable wedge. A sample composed only of safe toy APIs would hide the context and permission frictions that make the wedge interesting.

Figure~\ref{fig:paired} depicts four arms. The \emph{stock product} receives an ordinary owner request and available native tools. The \emph{expert-nudged product} receives the same task but permits a skilled operator to reframe, decompose, move context, and guide it to completion. The \emph{deterministic integration} uses conventional software where a stable API or rule-based workflow can do the job. The \emph{aftermarket harness} uses the same underlying model version as the stock arm but adds the specified task contract, context, control, and verification. All arms face the same terminal rule and the same legitimate authority boundary. The deterministic arm is a serious competitor, not a straw man.

\begin{figure}[tbp]
\centering\begin{tikzpicture}[x=1cm,y=1cm,>=Latex]
\fill[ink] (0,4.4) rectangle (12.7,5.0);
\node[anchor=west,text=white,font=\small\bfseries] at (.22,4.7) {Same held-out tasks and acceptance rule; different intervention};
\fill[ink!7] (0,0.65) rectangle (2.25,4.1);
\node[text=ink,font=\small\bfseries,align=center] at (1.12,2.6) {Held-out\\task pool};
\node[text=ink!70,font=\scriptsize,align=center] at (1.12,1.8) {Same starting state\\and consent};
\fill[ink!7] (10.55,.65) rectangle (12.7,4.1);
\node[text=ink,font=\small\bfseries,align=center] at (11.63,2.6) {Independent\\verifier};
\node[text=ink!70,font=\scriptsize,align=center] at (11.63,1.8) {Same terminal\\condition};
\def\testrow#1#2#3#4{
  \fill[#3!13] (3.05,#1-.31) rectangle (9.75,#1+.31);
  \draw[#3,line width=.9pt] (3.05,#1-.31)--(3.05,#1+.31);
  \node[anchor=west,text=ink,font=\small\bfseries] at (3.25,#1) {#2};
  \node[anchor=east,text=ink!70,font=\scriptsize] at (9.5,#1) {#4};
  \draw[->,ink!50,line width=.55pt] (2.25,#1)--(3.0,#1);
  \draw[->,ink!50,line width=.55pt] (9.80,#1)--(10.5,#1);
}
\testrow{3.65}{Stock consumer product}{native}{No extra intervention}
\testrow{2.83}{Expert nudging}{ochre}{Prompting and manual transfer}
\testrow{2.01}{Deterministic integration}{ink}{Conventional automation}
\testrow{1.19}{Aftermarket harness}{teal}{Productized context and control}
\node[anchor=west,text=ink!75,font=\scriptsize] at (.05,.2) {Record: accepted outcome $\cdot$ user-active minutes $\cdot$ latency $\cdot$ full cost $\cdot$ failure cause $\cdot$ repeatability};
\end{tikzpicture}

\caption{Proposed paired evaluation. Each arm receives the same held-out task definition and is judged by the same independent acceptance rule. The figure specifies a protocol and contains no results. Budget-matched and best-effort comparisons should be reported separately.}
\label{fig:paired}
\end{figure}

The stock and harness arms estimate a product-level completion difference. The expert-nudged and harness arms estimate whether the specialist has converted expert intervention into a usable product. The deterministic arm tests whether a language model was needed at all. A second design, an ablation within the harness, can isolate causal components: add structured task definition, then context retrieval, permitted tool access, durable state, exception handling, and independent verification one at a time. Because components interact, the order should be rotated or a factorial subset used where feasible. ``The harness helped'' is a start; knowing which element helped is the research result.

Each episode should be instrumented at the point of stop. Record the product version, model version if available, tools visible to each arm, account permissions, prompt history, elapsed time, user-active minutes, remote actions, and verifier evidence. Predeclare the primary cause of noncompletion and allow secondary labels. A decline is not coded as over-refusal merely because a later attempt succeeded; the new attempt may have gained a tool, permission, or context absent before. Conversely, a model that refuses a clearly benign task under the same conditions may reveal calibration trouble. XSTest and OR-Bench motivate careful benign-versus-prohibited labeling, while remaining distinct from this end-to-end commercial task suite \citep{xstest,orbench}.

The primary endpoint is \emph{verified authorized completion by the deadline}. A candidate answer without an external state change is partial. A platform response that has not yet settled is pending. A side effect whose result is ambiguous after a timeout is unknown. A correctly blocked request is a correct boundary outcome but not counted as completed work. Report these states separately. Secondary measures include accepted quality, human active minutes, number of nudges, latency, model and tool cost, recovery burden, duplicate effects, and repeated-run consistency. The $\mathrm{pass}^{k}$ formulation in $\tau$-bench emphasizes consistent success across repeated trials rather than a lucky success in one of several tries \citep{taubench}; a workflow study should similarly report variability and task-level failure patterns.

External verification matters because a harness can look effective by lowering its own standard. The acceptance rule should be written before running any arm and checked independently of the agent that performed the work. Where an API provides a durable record, use it. Where a person must judge an artifact, use blinded review and report disagreement. If a task cannot be verified without trusting the product's own summary, it should not be a headline completion result. This discipline is consistent with the functional end-state emphasis of WebArena and the complex digital-work tasks in TheAgentCompany \citep{webarena,agentcompany}. Those environments provide methodological precedents, not direct estimates of 2026 consumer-product refusal.

Holdouts should span task templates, organizations, account configurations, services, and time. A harness tuned to one website may pass familiar tasks while failing when an interface changes. A task suite should include new goals and routine exceptions, not only new wording for the same action. Report the evaluation budget, any human coaching allowed in each arm, and whether a run was stopped by cost, time, policy, or technical failure. \citet{agentsmatter} show why costs, holdouts, and reproducibility belong beside accuracy; leaving them out can make elaborate agents look better than cheaper alternatives.

There are two fair budget views. In a \emph{budget-matched} test, each arm gets the same maximum monetary and human-active allowance. In a \emph{best-effort} test, each arm can spend what a customer might realistically tolerate; the report then displays cost per accepted result rather than only success. A harness that wins under an unlimited budget but loses under customer constraints may be a useful research tool and a poor product. Both views should be published so the commercial argument cannot hide behind one favorable operating point.

Finally, an evaluator should report confidence intervals and case-level traces rather than a single victory number. The relevant cohort is likely heterogeneous: a harness may excel at multi-service tasks and add little to stable API tasks; human nudging may suffice for rare one-off work; a native provider upgrade may close a formerly large gap. These differences are exactly what a specialist needs to know before claiming a mass-market opportunity.
